\documentclass[11pt,a4paper]{tubaf-report}
\usepackage[sans]{tubaf-fonts}
\usepackage[english]{babel}
\usepackage{graphicx}
\graphicspath{{report/}{./}}

\usepackage[margin=2.4cm]{geometry}
\usepackage{amsmath,amssymb,bm}
\usepackage{booktabs,longtable,tabularx}
\usepackage{enumitem}
\usepackage{listings}
\usepackage{microtype}
\usepackage{tikz}
\usepackage{xcolor}
\usepackage[colorlinks=true,linkcolor=black,citecolor=black,urlcolor=blue!50!black]{hyperref}

\definecolor{codebg}{RGB}{246,246,246}
\lstset{
  basicstyle=\ttfamily\footnotesize,
  backgroundcolor=\color{codebg},
  frame=single,
  framerule=0.2pt,
  xleftmargin=5pt,
  xrightmargin=5pt,
  breaklines=true,
  columns=fullflexible,
  keepspaces=true,
  showstringspaces=false
}

\newcommand{\ud}{\,\mathrm{d}}
\newcommand{\vect}[1]{\bm{#1}}
\newcommand{\ten}[1]{\bm{#1}}
\newcommand{\B}{\mathbf{B}}
\newcommand{\N}{\mathbf{N}}
\newcommand{\trans}{\mathsf{T}}
\newcommand{\MPa}{\,\mathrm{MPa}}
\newcommand{\GPa}{\,\mathrm{GPa}}
\newcommand{\kPa}{\,\mathrm{kPa}}
\newcommand{\degC}{\,^\circ\mathrm{C}}

\title[OGS Training Course -- Benchmark Exercises]{Coupled Processes in Geological Media with OpenGeoSys}
\subtitle{Course exercises: tunnel excavation (Kirsch benchmark),
  fluid injection into fractured rock (LIE vs.\ EFPM),
  and thermo-hydro-mechanics of a heated tunnel (TRM)
}
\author[Ghasabeh and Nagel]{Mehran Ghasabeh \and Thomas Nagel}
\date{October 2026}
\subject{OpenGeoSys Training Course}
\publishers{Institute of Geotechnics, TU Bergakademie Freiberg}
\titlefoot{Joint OpenGeoSys Training Course -- UFZ, KIGAM, Zhejiang University, and TUBAF}
\setlength{\logoheadsize}{2.5cm}
\makeatletter
\renewcommand{\tubaf@title@logo}{\insertlogo{5.5cm}}
\makeatother
\logo{\includegraphics{UFZ_KIGAM_ZJU}}

\begin{document}
\maketitle
\tableofcontents
\clearpage

\chapter{Introduction}
\label{ch:intro}

This report documents the three benchmark exercises of the joint OpenGeoSys
training course. All exercises use the open-source finite-element simulator
OpenGeoSys (OGS)~\cite{bilke2019,kolditz2012,ogs-docs}, version 6.5.9,
installed as a Python wheel. The exercises increase in physical complexity:

\begin{enumerate}[nosep]
  \item \textbf{Tunnel excavation -- Kirsch benchmark}
        (\texttt{Tunnel\_Excavation\_Kirsch/}): linear-elastic excavation of a
        circular tunnel with the \texttt{ReleaseNodalForce} boundary
        condition, verified against the analytical Kirsch
        solution~\cite{kirsch1898}.
  \item \textbf{Fluid injection into fractured rock}
        (\texttt{LIE\_EFPM/}): hydro-mechanically coupled injection into a
        fractured, nearly impermeable matrix, comparing lower-dimensional
        interface elements (LIE) with the embedded fracture permeability
        model (EFPM), verified against analytical and one-dimensional
        finite-difference references.
  \item \textbf{Heated tunnel in Opalinus Clay}
        (\texttt{TRM/}): coupled thermo--hydro--mechanical response of a
        bentonite-buffered, heat-emitting tunnel modelled with the
        \path{THERMO_RICHARDS_MECHANICS} process, following the heater
        benchmark of Pitz et al.~\cite{pitz2023} (DECOVALEX-2023 Task~C), in a
        fully saturated and a partially saturated configuration.
\end{enumerate}

\section{Common workflow}

All three exercises follow the same reproducible structure:
\begin{itemize}[nosep]
  \item A single Jupyter notebook per problem ---
        \path{01_tunnel_excavation_kirsch.ipynb},
        \path{02_lie_efpm_injection.ipynb}, and
        \path{03_trm_heated_tunnel.ipynb} --- is the
        only place where parameter values are defined. It drives the
        complete workflow:
        meshing, project-file generation, OGS execution, post-processing, and
        verification.
  \item A mesh module (\texttt{mesh\_*.py}, class \texttt{MeshGenerator})
        builds structured Gmsh meshes, converts them to OGS \texttt{vtu}
        format, and extracts the boundary meshes.
  \item A project module (\texttt{prj\_*.py}, class \texttt{PrjGenerator} or
        \texttt{ProjectGenerator}) assembles the OGS project file with the
        \texttt{ogstools} \texttt{Project} API. All physical and numerical
        settings enter as required keyword arguments; the modules contain no
        hidden defaults.
\end{itemize}
This separation---\emph{values in the notebook, structure in small Python
modules, verification against analytical or published references}---is
intended as a reusable template for the participants' own models.

\section{Installation}
\label{sec:installation}

The course was developed and tested on Linux with Python~3.12; macOS and
Windows work analogously. All software installs into one virtual environment;
no compilation of OGS is required:
\begin{lstlisting}[language=bash]
python3 -m venv ~/.venvs/ogs312
source ~/.venvs/ogs312/bin/activate
pip install --upgrade pip
pip install ogs==6.5.9 ogstools==0.8.2 gmsh pyvista jupyter
python -m ipykernel install --user --name ogs312 --display-name "Python (ogs312)"
\end{lstlisting}
Open any course notebook in Jupyter and select the \texttt{Python (ogs312)}
kernel. Running all cells regenerates meshes and project files, runs OGS, and
produces the figures. Separate step-by-step installation guides for Windows,
macOS, and Ubuntu are provided alongside this report.

{\sloppy
One known limitation is relevant for the second exercise: in released OGS
versions, including the 6.5.9 wheel, the process
\path{HYDRO_MECHANICS_WITH_LIE} supports only a \emph{single}
displacement jump. The
double-fracture case is therefore run with the EFPM approach; double-fracture
LIE simulations require a custom build from the development branch
\path{LIE_HM_mutipleDisplacementJump} and are not part of the course.\par}

%=====================================================================
\chapter{Problem 1: Tunnel excavation and the Kirsch benchmark}
\label{ch:kirsch}

\section{Purpose}

Excavation removes both material and the tractions that this material exerted
on the surrounding rock. The OGS \texttt{ReleaseNodalForce} boundary condition
represents this unloading while retaining a mesh that already contains the
opening: the model begins from an initially stressed, undeformed state and
progressively withdraws the equivalent nodal forces on the excavation
boundary. After working through this exercise a participant should be able to
\begin{itemize}[nosep]
  \item distinguish an initial-stress excavation analysis from a model loaded
        from an unstressed reference configuration;
  \item derive the non-equilibrium initial residual and the equivalent nodal
        forces on the excavation wall;
  \item construct the plane-strain initial stress consistently; and
  \item verify OGS results against the Kirsch solution without confusing
        discretization error with nodal stress-extrapolation artifacts.
\end{itemize}
A dedicated companion report with the complete derivation is provided in the
folder \path{Tunnel_Excavation_Kirsch/report/}; this chapter
summarizes it.

\section{Benchmark definition}

A quarter model of a circular opening is used; the origin of the polar system
is the tunnel centre and $\theta$ is measured counter-clockwise from the
positive horizontal axis. Gravity is disabled; the in-situ stress is
prescribed directly.

\begin{center}
\begin{tabular}{@{}lll@{}}
\toprule
Quantity & Symbol & Value \\
\midrule
Tunnel radius & $a$ & $6.5\,\mathrm{m}$ \\
Quarter-domain extent & $L$ & $70\,\mathrm{m}$ \\
Young's modulus & $E$ & $5.0\GPa$ \\
Poisson's ratio & $\nu$ & $0.30$ \\
Vertical compressive stress magnitude & $P$ & $20\MPa$ \\
Excavation duration & $t_{\mathrm e}$ & $1.728\times10^5\,\mathrm{s}=2\,\mathrm d$ \\
Simulation end time & $t_{\mathrm{end}}$ & $3.456\times10^5\,\mathrm{s}=4\,\mathrm d$ \\
\bottomrule
\end{tabular}
\end{center}

OGS uses a tension-positive stress convention, so the applied far-field
stress is $\sigma_\infty=-P=-20\MPa$. Symmetry conditions are $u_x=0$ on the
left and $u_y=0$ on the bottom boundary; a vertical traction
$\bar t_y=-P$ acts on the top; the right boundary is traction-free,
consistent with the remote uniaxial stress state. The finite outer boundary
lies at $L/a\approx10.8$, so small differences from the infinite-domain
analytical solution are expected near the outer boundary.

\subsection{Plane-strain initial stress}
\label{sec:plane-strain}

The two-dimensional small-deformation process stores the stress vector as
$\vect{\sigma}=[\sigma_{xx},\sigma_{yy},\sigma_{zz},\sigma_{xy}]^{\trans}$.
Plane strain requires $\varepsilon_{zz}=0$, and isotropic Hooke's law then
yields the closure
\begin{equation}
  \boxed{\sigma_{zz}=\nu\left(\sigma_{xx}+\sigma_{yy}\right).}
  \label{eq:plane-strain-closure}
\end{equation}
For the intact uniaxial initial state, $\sigma_{0,xx}=0$ and
$\sigma_{0,yy}=-P$, hence
\begin{equation}
  \sigma_{0,zz}=-\nu P=-6\MPa,
  \qquad
  \vect{\sigma}_0=[0,-20,-6,0]^{\trans}\MPa.
  \label{eq:initial-stress}
\end{equation}
Omitting $\sigma_{0,zz}$ does not change the in-plane equilibrium equations,
but it shifts the computed $\sigma_{zz}$ by $6\MPa$ and therefore corrupts
comparisons of the full stress tensor, the mean stress, and the octahedral
shear stress.

\section{Analytical Kirsch solution}
\label{sec:kirsch-analytical}

For a traction-free circular opening of radius $a$ in an infinite isotropic
elastic medium under the remote vertical stress $\sigma_\infty$, the in-plane
stresses~\cite{kirsch1898} are
\begin{align}
\sigma_{rr}(r,\theta)
 &= \frac{\sigma_\infty}{2}\left(1-\frac{a^2}{r^2}\right)
  - \frac{\sigma_\infty}{2}
    \left(1-\frac{4a^2}{r^2}+\frac{3a^4}{r^4}\right)\cos 2\theta,
  \label{eq:kirsch-rr}\\
\sigma_{\theta\theta}(r,\theta)
 &= \frac{\sigma_\infty}{2}\left(1+\frac{a^2}{r^2}\right)
  + \frac{\sigma_\infty}{2}
    \left(1+\frac{3a^4}{r^4}\right)\cos 2\theta,
  \label{eq:kirsch-tt}\\
\sigma_{r\theta}(r,\theta)
 &= \frac{\sigma_\infty}{2}
    \left(1+\frac{2a^2}{r^2}-\frac{3a^4}{r^4}\right)\sin 2\theta.
  \label{eq:kirsch-rt}
\end{align}
At $r=a$: $\sigma_{rr}=\sigma_{r\theta}=0$ (traction-free wall),
$\sigma_{\theta\theta}=3\sigma_\infty=-60\MPa$ at the sidewall
($\theta=0$), and $\sigma_{\theta\theta}=-\sigma_\infty=+20\MPa$ at the
crown ($\theta=90^\circ$). The compatible out-of-plane stress follows from
Eq.~\eqref{eq:plane-strain-closure} applied to the rotation-invariant
in-plane trace:
\begin{equation}
  \sigma_{zz}(r,\theta)
  =\nu\left[\sigma_{rr}(r,\theta)+\sigma_{\theta\theta}(r,\theta)\right].
  \label{eq:kirsch-zz}
\end{equation}

\section{Release-nodal-force formulation}

For quasi-static small deformation, the discrete residual of the weak
equilibrium form reads
\begin{equation}
  \vect{r}(\hat{\vect{u}},t)
  =\vect{F}^{\mathrm{int}}(\hat{\vect{u}})
   -\vect{F}^{\mathrm{ext}}(t)=\vect{0},
  \qquad
  \vect{F}^{\mathrm{int}}
  =\int_\Omega \B^{\trans}\vect{\sigma}\ud\Omega,
  \qquad
  \ten{\sigma}=\ten{\sigma}_0+\mathbb{C}:\ten{\varepsilon}(\vect{u}).
  \label{eq:discrete-residual}
\end{equation}
At $t=0$ the mesh already contains the opening, carries the intact in-situ
stress $\ten{\sigma}_0$, and has zero displacement. This state is not in
equilibrium; OGS evaluates the \emph{non-equilibrium initial residual} once,
\begin{equation}
  \vect{r}^{\mathrm{neq}}
  =\int_\Omega\B^{\trans}\vect{\sigma}_0\ud\Omega
   -\vect{F}^{\mathrm{ext}}(0).
  \label{eq:rneq}
\end{equation}
Integration by parts shows that for a uniform (divergence-free) initial
stress, the free-degree entries of $\vect{r}^{\mathrm{neq}}$ live on the
excavation boundary $\Gamma_{\mathrm{exc}}$; the consistent equivalent nodal
force at wall node $i$ is
\begin{equation}
  \vect{F}_{0,i}
  =\int_{\Gamma_{\mathrm{exc}}}
    N_i^{\trans}\ten{\sigma}_0\vect{n}\ud\Gamma .
  \label{eq:F0}
\end{equation}
For the uniaxial initial stress this support traction has normal
component $t_n=-P\sin^2\theta$ and tangential component
$t_\theta=P\sin\theta\cos\theta$; it is a spatially varying traction field,
\emph{not} a uniform pressure (that interpretation is valid only for an
isotropic initial stress $\ten{\sigma}_0=-p_0\mathbf{I}$).

OGS subtracts the stored initial residual in every nonlinear iteration
(\emph{residual compensation}), so that $\vect{u}=\vect{0}$ satisfies
equilibrium at $t=0$; the compensation represents the support of the
not-yet-excavated core. The boundary condition then withdraws this support
with a user-defined decay function $f(t)$, $f(0)=1$, yielding the effective
wall residual
\begin{equation}
  \boxed{\tilde{\vect{r}}_i
  =\vect{F}^{\mathrm{int}}_i
   -\vect{F}^{\mathrm{ext}}_i
   -f(t)\vect{F}_{0,i}=\vect{0}.}
  \label{eq:effective-residual}
\end{equation}
The notebook uses the linear decay
\begin{equation}
  f(t)=
  \begin{cases}
    1-t/t_{\mathrm e}, & 0\le t\le t_{\mathrm e},\\[4pt]
    0, & t>t_{\mathrm e},
  \end{cases}
  \label{eq:decay}
\end{equation}
so the wall is in-situ supported at $t=0$, partially unloaded for
$0<t<t_{\mathrm e}$, and traction-free for $t\ge t_{\mathrm e}$. Because the
material is linear elastic, the excavation-induced response scales linearly
with the released fraction $1-f(t)$.

\begin{center}
\begin{tikzpicture}[scale=0.95]
  \draw[->] (0,0) -- (7,0) node[right] {$t$};
  \draw[->] (0,0) -- (0,3) node[above] {$f(t)$};
  \draw[very thick,blue!60!black] (0,2.3) -- (4.5,0) -- (6.5,0);
  \draw[dashed] (4.5,0) -- (4.5,2.3);
  \draw[dashed] (0,2.3) -- (4.5,2.3);
  \node[left] at (0,2.3) {$1$};
  \node[below] at (4.5,0) {$t_{\mathrm e}$};
  \node[blue!60!black] at (2.8,1.6) {$1-t/t_{\mathrm e}$};
\end{tikzpicture}
\end{center}

\section{OGS realization}

The process variable requests the residual compensation and the wall boundary
condition:
\begin{lstlisting}[language=XML]
<process_variable>
  <name>displacement</name>
  <components>2</components>
  <order>2</order>
  <initial_condition>initial_displacement</initial_condition>
  <compensate_non_equilibrium_initial_residuum>true
  </compensate_non_equilibrium_initial_residuum>
  <boundary_conditions>
    <boundary_condition>
      <mesh>arc</mesh>
      <type>ReleaseNodalForce</type>
      <time_decay_parameter>decay_function</time_decay_parameter>
    </boundary_condition>
  </boundary_conditions>
</process_variable>
\end{lstlisting}
The mechanics process references the initial-stress parameter
\texttt{sigma0} through its \path{initial_stress} tag, with the
plane-strain components of Eq.~\eqref{eq:initial-stress}, and the decay
function of Eq.~\eqref{eq:decay} enters as
\begin{lstlisting}[language=XML]
<parameter>
  <name>decay_function</name>
  <type>Function</type>
  <expression>if (t>172800) 0.0; else 1-t/172800</expression>
</parameter>
\end{lstlisting}

Two analyses are compared: the \emph{release} formulation (initial stress,
\texttt{ReleaseNodalForce} wall, compensation enabled) and a \emph{standard}
formulation (zero initial stress, natural traction-free wall, load applied in
one step from the unstressed configuration). The final total stresses must
agree; the displacement fields need not, because they refer to different
reference configurations.

\section{Verification results}

The executed notebook (OGS 6.5.9) runs five simulations; all return code zero
and all final stress checks are asserted programmatically.

\begin{center}
\begin{tabularx}{\textwidth}{@{}Xrrr@{}}
\toprule
Measure & Analytical & Quadratic & Linear \\
\midrule
$\sigma_{rr}(a,0^\circ)$ / MPa & $0$ & $-0.133$ & $-2.620$ \\
$\sigma_{\theta\theta}(a,0^\circ)$ / MPa & $-60$ & $-61.581$ & $-62.476$ \\
Release--standard maximum stress difference / MPa
 & $0$ & $9.9\times10^{-11}$ & $4.5\times10^{-12}$ \\
Initial release-model displacement / m & $0$ & $0$ & $0$ \\
\bottomrule
\end{tabularx}
\end{center}

The release and standard stress fields agree to numerical precision. The
nonzero extrapolated $\sigma_{rr}$ at the wall is not a failure of the
traction boundary condition: OGS computes stress at integration points and
extrapolates to nodes for VTK output, and this postprocessing step, together
with mesh resolution and element order, creates the wall-node discrepancy.
Quadratic elements improve the wall value substantially on the present mesh.

An element-deactivation variant fills the opening with a quarter-disc of
inactive elements (material ID~1). OGS assigns zero integration-point stress
to inactive cells, so nodal extrapolation at the shared wall nodes mixes
physical rock stresses with zeros; the notebook therefore excludes the shared
wall node from the checks. Away from that interface, the deactivation and
hole-only in-plane stresses agree within $10^{-8}\MPa$. For wall
verification one should use integration-point or cell-based stresses, or a
sample point demonstrably inside the active rock.

The notebook additionally demonstrates polar postprocessing
(\path{ot.variables.stress.to_polar}) and maps the mean stress
$\sigma_{\mathrm m}=\tfrac13\operatorname{tr}\ten{\sigma}$ and octahedral
shear stress $\tau_{\mathrm{oct}}=\sqrt{2J_2/3}$ --- the invariant pair
relevant for pressure-dependent failure criteria: excavation simultaneously
concentrates shear and removes confinement at the wall.

%=====================================================================
\chapter{Problem 2: Fluid injection into fractured rock (LIE vs.\ EFPM)}
\label{ch:lie}

\section{Purpose and background}

The second exercise addresses hydro-mechanically coupled processes in
fault-controlled reservoirs, motivated by deep geothermal systems in
fractured carbonates. A pre-existing fracture in a nearly impermeable rock
matrix is pressurised by fluid injection and subsequently depressurised by
production. Two modelling strategies are benchmarked against each other and
against analytical references:
\begin{itemize}[nosep]
  \item \textbf{LIE} --- the fracture is a lower-dimensional interface
        element with an explicit displacement-jump unknown
        (process \path{HYDRO_MECHANICS_WITH_LIE}); fracture flow follows
        the cubic law with pressure-dependent aperture.
  \item \textbf{EFPM} --- the fracture is a thin band of standard continuum
        elements (process \path{HYDRO_MECHANICS}) whose permeability and
        stiffness are \emph{equivalent} upscaled properties: the medium
        property
        \texttt{Orthotropic\allowbreak Embedded\allowbreak Fracture\allowbreak \allowbreak Permeability}
        and a Gerrard-type equivalent modulus~\cite{gerrard1982}.
\end{itemize}
Learning objectives: discrete versus equivalent-continuum fracture
representations, the cubic law, hydro-mechanical aperture coupling, and a
verification methodology combining closed-form and numerical references.

\section{Model definition}

\subsection{Geometry and mesh}

The domain is a plane-strain rock-matrix block of length
$12.5\,\mathrm m$ ($x\in[0.1,12.6]\,\mathrm m$) and height $1\,\mathrm m$. A
single horizontal fracture runs along $y=0.5\,\mathrm m$ from the left to
the right edge. In the LIE model it is a line of 1D interface elements (34
elements embedded in 1292 quadratic quads); in the EFPM model it is a band of
thickness $s=1\,\mathrm{mm}$, one element thick (39 band elements within
2301 quads). The double-fracture variant places fractures at
$y=0.3\,\mathrm m$ and $y=0.7\,\mathrm m$. All meshes are structured,
transfinite, graded towards the fracture, and converted to quadratic
elements; displacement is quadratic, pressure linear (Taylor--Hood).

\subsection{Fracture flow and the two representations}

Fracture flow obeys the \emph{cubic law}~\cite{witherspoon1980}: a fracture
of hydraulic aperture $b$ has permeability and transmissivity
\begin{equation}
  k_{\mathrm F}=\frac{b^2}{12},
  \qquad
  T_{\mathrm{LIE}}=k_{\mathrm F}\,b=\frac{b^3}{12}.
  \label{eq:cubic}
\end{equation}
In the LIE model the aperture evolves with the normal displacement jump; in
linearized form $b=b_0+\Delta p/K_n$, and the fracture storage combines
fluid storage and fracture compliance,
$\beta=b\,S_{\mathrm F}+1/K_n$.

In the EFPM the band of width $s$ carries the upscaled permeability
\begin{equation}
  \ten k = k_{\mathrm M}\ten I
         + \frac{b}{a}\,\frac{b^2}{12}
           \left(\ten I-\vect n\otimes\vect n\right),
  \qquad
  b = b_0 + a\,\langle\varepsilon_n-\varepsilon_0\rangle ,
  \label{eq:efpm-k}
\end{equation}
where $a$ is the mean fracture spacing, $\vect n$ the fracture normal, and
$\varepsilon_n$ the strain normal to the fracture. Choosing $s=a$ makes the
band transmissivity $T_{\mathrm{EFPM}}=k\,s$ equal to
Eq.~\eqref{eq:cubic}. With $b_0=0$, the initial aperture
$b(t{=}0)=10\,\mu\mathrm m$ is encoded by the threshold strain
$\varepsilon_0=-b(t{=}0)/a=-10^{-2}$. The band's normal compliance mimics the
fracture stiffness through the Gerrard equivalent
modulus~\cite{gerrard1982},
\begin{equation}
  \frac{1}{E_{\mathrm{RM}}}
  =\frac{1}{E_{\mathrm M}}+\frac{1}{K_n\,a}
  \quad\Rightarrow\quad
  E_{\mathrm{RM}}=9.98\times10^{7}\,\mathrm{Pa}.
  \label{eq:gerrard}
\end{equation}

\subsection{Parameters}

\begin{center}
\begin{tabular}{@{}lll@{}}
\toprule
Property & Symbol & Value \\
\midrule
\multicolumn{3}{@{}l@{}}{\emph{Rock matrix}} \\
Young's modulus & $E$ & $60\GPa$ \\
Poisson's ratio & $\nu$ & $0$ \\
Solid density & $\varrho_{\mathrm s}$ & $2716\,\mathrm{kg/m^3}$ \\
Porosity & $\phi$ & $0.001$ \\
Permeability & $k_{\mathrm M}$ & $10^{-21}\,\mathrm{m^2}$ \\
Biot coefficient & $\alpha_{\mathrm B}$ & $1$ \\
Storage & $S$ & $10^{-10}\,\mathrm{1/Pa}$ \\
\multicolumn{3}{@{}l@{}}{\emph{Fluid}} \\
Density & $\varrho_{\mathrm{LR}}$ & $1000\,\mathrm{kg/m^3}$ \\
Viscosity & $\mu$ & $10^{-3}\,\mathrm{Pa\,s}$ \\
Compressibility & $\kappa_{\mathrm f}$ & $10^{-7}\,\mathrm{1/Pa}$ \\
\multicolumn{3}{@{}l@{}}{\emph{Fracture (LIE)}} \\
Normal and shear stiffness & $K_n=K_s$ & $10^{11}\,\mathrm{Pa/m}$ \\
Initial aperture & $b_0$ & $10^{-5}\,\mathrm m$ \\
Fracture storage & $S_{\mathrm F}$ & $10^{-10}\,\mathrm{1/Pa}$ \\
Permeability model & & \texttt{CubicLawPermeability} \\
\multicolumn{3}{@{}l@{}}{\emph{Fracture band (EFPM)}} \\
Band thickness, fracture spacing & $s=a$ & $10^{-3}\,\mathrm m$ \\
Threshold strain & $\varepsilon_0$ & $-10^{-2}$ \\
Band Young's modulus & $E_{\mathrm{RM}}$ & $9.98\times10^{7}\,\mathrm{Pa}$ \\
Fracture normal & $\vect n$ & $(0,1,0)$ \\
\bottomrule
\end{tabular}
\end{center}

\subsection{Initial conditions, boundary conditions, and injection schedule}

The LIE model works with absolute quantities: initial and far-field pore
pressure $p_0=11\MPa$ (also a Dirichlet condition at the outlet point on
the right edge), vertical total load $\sigma_{yy}=-50\MPa$ on the top
boundary, and an initial effective stress of
$\sigma'_{yy}=-39\MPa$ in the matrix and on the fracture. The EFPM model is
formulated in \emph{relative} quantities (zero initial pressure, load, and
stress); its results are compared after adding $p_0$. Displacement
conditions are $u_y=0$ on the bottom and $u_x=0$ on the left and right
boundaries.

Fluid is injected at the left fracture end at the mass rate
$\dot m=2\times10^{-5}\,\mathrm{kg/(s\,m)}$, i.e.\ a volumetric rate
$Q=2\times10^{-8}\,\mathrm{m^3/(s\,m)}$: in the LIE model as a point source
at the fracture inlet node, in the EFPM model as the equivalent flux
$Q/s=2\times10^{-2}\,\mathrm{m/s}$ on the $1\,\mathrm{mm}$ inlet edge of the
band. The schedule ramps the rate up over $0$--$100\,\mathrm s$, holds it
constant until $1000\,\mathrm s$, reverses it linearly over
$1000$--$1100\,\mathrm s$ (switch from injection to production), and holds
the reversed rate until the end time $t_{\mathrm{end}}=1500\,\mathrm s$.

\subsection{Numerical settings}

Backward Euler with iteration-number-based time stepping (initial
$\Delta t=10\,\mathrm s$, $\Delta t\in[0.1,20]\,\mathrm s$); damped Newton
(damping factor $0.8$, maximum 50 iterations) with per-component
$\Delta x$ criteria; Eigen \texttt{SparseLU} with scaling as the linear
solver. All three course cases run with the stock OGS~6.5.9 wheel:
single-fracture LIE in $83\,\mathrm s$ (76 steps), single-fracture EFPM in
$357\,\mathrm s$ (95 steps), double-fracture EFPM in about $3$ minutes
(76 steps), with no rejected steps. The systems are ill-conditioned (the
$K_n=10^{11}\,\mathrm{Pa/m}$ penalty-like fracture stiffness), which is why
a direct solver is used; on builds with Intel MKL, \texttt{PardisoLU} solves
the same systems roughly $2.8\times$ faster.

\section{Verification and results}

\subsection{Analytical and 1D references}

Two independent references are built in the notebook for the single-fracture
case:
\begin{enumerate}[nosep]
  \item a \emph{linearized closed-form solution} of the fracture diffusion
        problem of Carslaw--Jaeger type~\cite{carslaw1959}
        (constant-flux half-line, $\operatorname{ierfc}$ kernel), with the
        $100\,\mathrm s$ injection ramp incorporated exactly via Duhamel's
        principle; and
  \item a \emph{nonlinear 1D finite-difference solution} of
        \begin{equation}
          \beta(b)\,\frac{\partial p}{\partial t}
          =\frac{\partial}{\partial x}\!\left(
          \frac{b^3}{12\mu}\frac{\partial p}{\partial x}\right),
          \qquad
          b=b_0+\frac{p-p_0}{K_n},
          \qquad
          \beta=b\,S_{\mathrm F}+\frac{1}{K_n},
          \label{eq:fd-ref}
        \end{equation}
        which retains the cubic-law nonlinearity and the ramped source
        $Q(t)$.
\end{enumerate}
Including the injection ramp is essential: without it both references lead
the simulations systematically. The initial transmissivity, storativity, and
hydraulic diffusivity are
$T_0=8.33\times10^{-14}\,\mathrm{m^3/(Pa\,s)}$,
$\beta_0=1.00\times10^{-11}\,\mathrm{m/Pa}$, and
$D_0=8.33\times10^{-3}\,\mathrm{m^2/s}$.

\subsection{Key results}

\begin{center}
\begin{tabular}{@{}lccc@{}}
\toprule
Quantity (single fracture) & LIE & EFPM & 1D reference \\
\midrule
Peak inlet pressure / MPa & $11.489$ & $11.480$ & --- \\
Inlet overpressure at $t=1000\,\mathrm s$ / kPa & $489.0$ & --- & $505.1$ \\
Inlet pressure at $t=1500\,\mathrm s$ / MPa & $10.808$ & $10.795$ & --- \\
Maximum $|u_y|$ at $t\approx1000\,\mathrm s$ / $\mu$m & $4.79$ & $4.69$ & --- \\
\bottomrule
\end{tabular}
\end{center}

\begin{itemize}[nosep]
  \item \textbf{LIE vs.\ EFPM.} Inlet and outlet pressure histories and
        pressure profiles along the fracture agree within about
        $0.1$--$0.2\,\%$ throughout injection and production; displacement
        fields agree within a few percent. The equivalent-continuum
        representation therefore reproduces the discrete-fracture response
        for this configuration.
  \item \textbf{LIE vs.\ analytical references.} The simulated inlet
        overpressure at the end of injection deviates by $3.3\,\%$ from the
        nonlinear 1D reference; the residual difference is the 2D elastic
        compliance of the matrix, which the 1D references do not contain.
  \item \textbf{Aperture coupling.} During injection the inlet aperture
        opens from $10$ to about $14.7\,\mu\mathrm m$
        (fracture permeability from $1.0\times10^{-11}$ to
        $1.8\times10^{-11}\,\mathrm{m^2}$, consistent with $b^2/12$ ---
        i.e.\ the transmissivity roughly triples), while the outlet aperture
        stays at $10\,\mu\mathrm m$; during production the inlet aperture
        closes below its initial value, to about $8.8\,\mu\mathrm m$ at
        $t=1500\,\mathrm s$. The pressure pulse is still propagating outward
        along the fracture while the inlet is already depressurising.
  \item \textbf{Double fracture (EFPM).} Both fractures respond nearly
        identically to each other and to the single-fracture case
        (peak inlet pressures within $0.01\MPa$); the maximum vertical
        displacement roughly doubles to about $9.4\,\mu\mathrm m$ because
        both fractures dilate.
\end{itemize}

As noted in Section~\ref{sec:installation}, released OGS versions support
only one displacement-jump variable, so the double-fracture configuration is
run with the EFPM approach in the course; a double-fracture LIE project file
is included for participants with access to the development branch.

%=====================================================================
\chapter{Problem 3: Heated tunnel in Opalinus Clay (TRM)}
\label{ch:trm}

\section{Purpose and background}

The third exercise models the near field of a heat-emitting tunnel --- a
plane-strain abstraction of the FE heater experiment at the Mont Terri rock
laboratory --- with the OGS \path{THERMO_RICHARDS_MECHANICS} (TRM)
process. It follows the systematic benchmark of Pitz et
al.~\cite{pitz2023} (DECOVALEX-2023 Task~C); the reference paper is included
in the exercise folder. Learning objectives:
\begin{itemize}[nosep]
  \item set up a coupled THM model under the Richards assumption;
  \item understand thermal pressurisation of pore water in a low-permeability
        host rock and its anisotropic drainage;
  \item model unsaturated behaviour: van~Genuchten retention and relative
        permeability, Bishop effective stress, vapour diffusion, latent heat,
        saturation-dependent thermal conductivity;
  \item use orthotropic materials with local coordinate systems; and
  \item work with mixed-order (Taylor--Hood) elements and interpolated
        output fields.
\end{itemize}

\section{Model definition}

\subsection{Geometry and mesh}

A quarter of a $50\,\mathrm m\times50\,\mathrm m$ domain is discretised
(quarter dimensions $25\,\mathrm m\times25\,\mathrm m$). A granular
bentonite buffer ring (material ID~0) extends from the heater surface at
$r_i=0.525\,\mathrm m$ to the tunnel wall at $R=1.24\,\mathrm m$; the
surrounding Opalinus Clay is material ID~1. The heater itself is not meshed;
it acts through a heat-flux boundary condition on the inner arc.

The mesh is a structured, transfinite quadrilateral Gmsh mesh (15 nodes per
block edge, 8 radial nodes across the buffer), converted to 8-node quadratic
quadrilaterals: 1176 cells (196 buffer, 980 clay) and 3683 nodes. The
displacement is interpolated with quadratic shape functions and pressure and
temperature with linear ones (Taylor--Hood). Boundary meshes
\texttt{left/right/top/bottom/arc\_inner/arc\_outer} are extracted
automatically.

\subsection{Process and governing equations}

The TRM process solves, in the domain $\Omega$ and for the unknowns
temperature $T$, liquid pressure $p_{\mathrm{LR}}$, and displacement
$\vect u$:
\begin{itemize}[nosep]
  \item the \textbf{energy balance} of the fluid--solid mixture with
        effective volumetric heat capacity
        $(\varrho c_p)_{\mathrm{eff}}$, heat conduction with the effective
        conductivity $\ten{\lambda}_{\mathrm{eff}}$, and advection by the
        Darcy flux (plus vapour-related latent-heat transport in the
        unsaturated configuration);
  \item the \textbf{liquid mass balance} under the Richards assumption (gas
        at constant atmospheric pressure), with storage from fluid and pore
        compressibility, saturation changes
        $\phi\,\partial S_{\mathrm L}/\partial t$, thermal expansion of the
        pore water, Darcy flow
        $ \vect q = -\dfrac{k_{\mathrm{rel}}(S_{\mathrm L})\,\ten k}{\mu(T)}
          \nabla p_{\mathrm{LR}}$,
        and, in the unsaturated configuration, diffusive water-vapour
        transport; and
  \item the \textbf{momentum balance}
        $\nabla\cdot\left(\ten\sigma' - \alpha_{\mathrm B}\,
        \chi(S_{\mathrm L})\,p_{\mathrm{LR}}\ten I\right)=\vect 0$ with the
        Bishop factor $\chi$ and thermo-elastic effective stress
        $\ten\sigma'=\mathbb{C}:\left(\ten\varepsilon
        -\alpha_{T}\Delta T\,\ten I\right)$.
\end{itemize}
Gravity is disabled. With $S_{\mathrm L}\equiv1$, $k_{\mathrm{rel}}\equiv1$,
and $\chi=S_{\mathrm L}$ the formulation reduces to single-phase saturated
THM --- this is the \emph{saturated} configuration (TRM-S
in~\cite{pitz2023}). The \emph{partially saturated} configuration activates
van~Genuchten saturation and relative permeability, a Bishop
saturation-cutoff model, water-vapour density and FEBEX-type vapour
diffusion with linear latent heat, mass lumping, and a temperature-dependent
IAPWS viscosity curve.

\subsection{Material parameters}

\begin{center}
\footnotesize
\begin{tabular}{@{}llll@{}}
\toprule
Property & Symbol & Bentonite buffer & Opalinus Clay \\
\midrule
Solid density / $\mathrm{kg\,m^{-3}}$ & $\varrho_{\mathrm s}$ & $2227.2$ & $2689.66$ \\
Solid specific heat / $\mathrm{J\,kg^{-1}\,K^{-1}}$ & $c_{p,\mathrm s}$ & $800$ & $995$ \\
Solid thermal conductivity / $\mathrm{W\,m^{-1}\,K^{-1}}$ & $\lambda_{\mathrm s}$ & $1.2$ &
  $\lambda_\perp=1.3$, $\lambda_\parallel=2.4$ \\
Solid thermal expansivity / $\mathrm{K^{-1}}$ & $\alpha_{T}$ & $3.0\times10^{-6}$ & $1.5\times10^{-5}$ \\
Intrinsic permeability / $\mathrm{m^2}$ & $k$ & $3.5\times10^{-20}$ &
  $k_\perp=10^{-20}$, $k_\parallel=5\times10^{-20}$ \\
Porosity & $\phi$ & $0.331$ & $0.13$ \\
Elasticity & & isotropic, $E=18\MPa$, $\nu=0.35$ &
  orthotropic, see text \\
Biot coefficient & $\alpha_{\mathrm B}$ & $1.0$ & $1.0$ \\
van Genuchten $p_{\mathrm b}$, $m$ & & $28.6\MPa$, $0.5$ & $20\MPa$, $0.6$ \\
\bottomrule
\end{tabular}
\end{center}

The clay is modelled as a linearly elastic orthotropic material with
$E_\perp=4\GPa$ and $E_\parallel=8\GPa$, Poisson's ratios
$(0.25,\,0.35,\,0.35)$, and shear moduli
$G=3.5\GPa$ and $8/(2\cdot1.35)\approx2.96\GPa$; its anisotropy
axes are set through an OGS local coordinate system
(\path{<local_coordinate_system>}), with bedding-parallel and
bedding-perpendicular directions as listed above. Two saturated bedding
orientations ($0^\circ$ and $90^\circ$) are generated from the same
parameter set by rotating the local basis and the permeability tensor.

The pore water uses $c_p=4184\,\mathrm{J/(kg\,K)}$,
$\lambda=0.63\,\mathrm{W/(m\,K)}$, viscosity
$\mu=10^{-3}\,\mathrm{Pa\,s}$, and a linear equation of state
$\varrho_{\mathrm{LR}}(p,T)$ with reference density
$1000\,\mathrm{kg/m^3}$, thermal expansivity $4\times10^{-4}\,\mathrm{1/K}$,
and compressibility $5\times10^{-10}\,\mathrm{1/Pa}$. The partially
saturated configuration replaces the constant viscosity by a 37-point IAPWS
$\mu(T)$ curve ($273.15$--$453.15\,\mathrm K$) and uses a buffer thermal
conductivity that increases linearly with saturation from
$0.35$ (dry) to $1.2\,\mathrm{W/(m\,K)}$ (saturated).

\subsection{Initial and boundary conditions}

\begin{center}
\begin{tabularx}{\textwidth}{@{}lXX@{}}
\toprule
Field & Saturated configuration & Partially saturated configuration \\
\midrule
$T_0$ & $288.15\,\mathrm K$ ($15\degC$) & $288.15\,\mathrm K$ \\
$p_0$ & $0\,\mathrm{Pa}$ (results shifted by $+2\MPa$ in the plots) &
  buffer $-123\MPa$ ($S_{\mathrm L}\approx0.23$, dry emplacement);
  clay $+2\MPa$ \\
$\vect u_0$, $\ten\sigma_0$ & zero & zero \\
\midrule
Displacement & \multicolumn{2}{p{0.72\textwidth}}{rollers on all four outer
  boundaries ($u_x=0$ left/right, $u_y=0$ top/bottom); heater surface
  (\texttt{arc\_inner}) fixed} \\
Temperature & \multicolumn{2}{p{0.72\textwidth}}{Neumann heat flux
  $q=89\,\mathrm{W/m^2}$ on \texttt{arc\_inner}, constant over the full five
  years} \\
Hydraulic & \multicolumn{2}{p{0.72\textwidth}}{all boundaries natural (no
  flow)} \\
\bottomrule
\end{tabularx}
\end{center}

Because the saturated response is linear in the pressure level, running it
with $p_0=0$ and shifting the plotted pressures by $+2\MPa$ reproduces the
paper's absolute values exactly.

\subsection{Numerical settings}

Backward Euler in time with iteration-number-based adaptive time stepping
(initial $\Delta t=0.1\,\mathrm s$, maximum
$\Delta t=1.576\times10^6\,\mathrm s$) up to
$t_{\mathrm{end}}=1.5768\times10^8\,\mathrm s$ (5 years). Newton's method
(maximum 20 iterations) with a per-component relative $\Delta x$ criterion of
$10^{-6}$ for $T$, $p$, $u_x$, $u_y$. The linear solver is Eigen's
\texttt{SparseLU} direct solver with equilibration
(\texttt{<scaling>true</scaling>}); iterative presets
(BiCGSTAB/GMRES with ILUT) are available in the notebook as named
alternatives. Because $T$ and $p$ are first-order fields on a second-order
mesh, the project outputs \texttt{temperature\_interpolated} and
\texttt{pressure\_interpolated}, which carry values on \emph{all} nodes of
the quadratic mesh; all probes and contour plots use these fields.

\section{Results}

All three runs (saturated $0^\circ$, saturated $90^\circ$, partially
saturated) complete with OGS 6.5.9 without any rejected time steps
(139/139/142 adaptive steps; about $80$--$150\,\mathrm s$ wall time each).
Representative values at the end of the five-year heating phase, probed on
the horizontal symmetry axis (saturated pressures include the $+2\MPa$
shift):

\begin{center}
\begin{tabular}{@{}lccc@{}}
\toprule
 & Heater surface & Tunnel wall & Clay \\
 & $r=0.525\,\mathrm m$ & $r=1.24\,\mathrm m$ & $r=3\,\mathrm m$ \\
\midrule
\multicolumn{4}{@{}l@{}}{\emph{Saturated ($90^\circ$ bedding case)}} \\
Temperature & $123.9\degC$ & $84.4\degC$ & $63.4\degC$ \\
Liquid pressure & $18.5\MPa$ & $18.4\MPa$ & $17.4\MPa$ \\[2pt]
\multicolumn{4}{@{}l@{}}{\emph{Partially saturated}} \\
Temperature & $161.7\degC$ & $79.8\degC$ & $54.4\degC$ \\
Liquid pressure & $-478\MPa$ & $-32.4\MPa$ & $+2.5\MPa$ \\
Saturation & $0.06$ & $0.52$ & $1.00$ \\
\bottomrule
\end{tabular}
\end{center}

\paragraph{Saturated configuration.}
Heating pressurises the pore water of the low-permeability clay to a peak of
about $16.5\MPa$ above the initial state (thermal pressurisation), which
then drains anisotropically along the bedding. The two bedding orientations
$0^\circ$ and $90^\circ$ produce stress, pressure, and displacement fields
that are exact $90^\circ$ rotations of each other (maximum displacement
$6.08\,\mathrm{mm}$ in both runs), a built-in consistency check of the local
coordinate system and tensor rotation handling.

\paragraph{Partially saturated configuration.}
The dry-emplaced buffer first desiccates further near the heater
(saturation drops from $0.23$ to about $0.06$, suction grows to several
hundred MPa) because the strong thermal gradient drives vapour outward,
while the buffer simultaneously resaturates from the host rock at the tunnel
wall. The clay pressurises to a peak of about $11\MPa$. The heater
temperature is markedly higher than in the saturated case
($\approx162\degC$ vs.\ $\approx124\degC$) because the dry buffer conducts
heat poorly; as the buffer resaturates its conductivity rises and the heater
temperature declines late in the simulation. This is the characteristic
behaviour reported in~\cite{pitz2023}; digitized reference curves from the
paper are archived in \texttt{TRM/paper\_data/} for further comparison
exercises.

\paragraph{Mesh sensitivity.}
A fine-mesh reference solution (5700 cells, archived in the folder
\path{TRM/_out_fine_reference/}) confirms the saturated
fields on the course mesh to within a fraction of a degree and
$0.1\MPa$. The steep dry-out front of the partially saturated case is more
mesh-sensitive: the fine mesh gives $151\degC$ and $-380\MPa$ at the heater
against $162\degC$ and $-478\MPa$ on the coarse course mesh. Participants
should treat near-heater suction values as mesh-dependent and refine locally
when they matter.

%=====================================================================
\chapter{Concluding remarks}

The three exercises cover the span from single-physics verification to
multi-field, multi-material applications:
\begin{itemize}[nosep]
  \item Problem~1 establishes the methodology: a formal derivation of the
        excavation boundary condition, exact agreement between two
        formulations, and a quantitative comparison with a classical
        analytical solution --- including the discipline of separating model
        error, discretization error, and postprocessing artifacts.
  \item Problem~2 contrasts two fracture representations inside one
        verification framework and shows that an equivalent-continuum model,
        when its transmissivity and compliance are upscaled consistently,
        matches a discrete interface-element model to within a fraction of a
        percent.
  \item Problem~3 scales this methodology to a coupled THM problem with two
        materials, anisotropy, and unsaturated constitutive behaviour,
        reproducing the published heater benchmark behaviour and including a
        mesh-sensitivity assessment.
\end{itemize}
All folders are self-contained and re-runnable from scratch with the single
Python environment of Section~\ref{sec:installation}; the executed course
versions were verified with OpenGeoSys~6.5.9.

\KOMAoptions{bibliography=totoc}
\begin{thebibliography}{9}
  \bibitem{kirsch1898}
        E. G. Kirsch, ``Die Theorie der Elastizit\"at und die Bed\"urfnisse
        der Festigkeitslehre,'' \emph{Zeitschrift des Vereines Deutscher
        Ingenieure}, vol.~42, pp.~797--807, 1898.
  \bibitem{pitz2023}
        M. Pitz et al., ``Non-isothermal consolidation: A systematic
        evaluation of two implementations based on multiphase and Richards
        equations,'' \emph{International Journal of Rock Mechanics and
        Mining Sciences}, vol.~170, 105534, 2023.
        \url{https://doi.org/10.1016/j.ijrmms.2023.105534}
  \bibitem{bilke2019}
        L. Bilke, B. Flemisch, T. Kalbacher, O. Kolditz, R. Helmig, and
        T. Nagel, ``Development of open-source porous media simulators:
        Principles and experiences,'' \emph{Transport in Porous Media},
        vol.~130, pp.~337--361, 2019.
  \bibitem{kolditz2012}
        O. Kolditz, U.-J. G\"orke, H. Shao, and W. Wang (eds.),
        \emph{Thermo-Hydro-Mechanical-Chemical Processes in Porous Media:
        Benchmarks and Examples}, Springer, Berlin, Heidelberg, 2012.
  \bibitem{witherspoon1980}
        P. A. Witherspoon, J. S. Y. Wang, K. Iwai, and J. E. Gale,
        ``Validity of cubic law for fluid flow in a deformable rock
        fracture,'' \emph{Water Resources Research}, vol.~16, no.~6,
        pp.~1016--1024, 1980.
  \bibitem{gerrard1982}
        C. M. Gerrard, ``Elastic models of rock masses having one, two and
        three sets of joints,'' \emph{International Journal of Rock
        Mechanics and Mining Sciences \& Geomechanics Abstracts}, vol.~19,
        pp.~15--23, 1982.
  \bibitem{carslaw1959}
        H. S. Carslaw and J. C. Jaeger, \emph{Conduction of Heat in Solids},
        2nd ed., Oxford University Press, 1959.
  \bibitem{ogs-docs}
        OpenGeoSys documentation, \url{https://www.opengeosys.org/docs/}.
\end{thebibliography}

\paragraph{Acknowledgement.}
The Kirsch excavation example builds on the OpenGeoSys
\texttt{ReleaseNodalForce} benchmark by Wenqing Wang; the TRM exercise
follows the DECOVALEX-2023 Task~C benchmark of Pitz et al.

\end{document}
